\section{Dataset and Task}

\subsection{Task Definition}

Each instance contains claim text $x_t$, an image $x_v$, the earliest observed claim time, a decision cutoff $\tau$, and an evidence collection $E_{\leq\tau}$. The output is $y \in \{\text{supported},\text{refuted},\text{insufficient evidence}\}$. Selective deferral is an operational decision applied after calibration, not a fourth class. The event is the independent unit for splitting and statistical analysis.

\subsection{Candidate Acquisition}

Corpus construction starts with English ClaimReview records about United States events from 2020--2026. GDELT supplies discovery metadata, and Common Crawl supplies archived page observations. The claim time is the earliest documented public appearance. The decision cutoff is the fact-check publication time or a recorded equivalent. The corpus retains cases with a recoverable image or archive, an unambiguous professional rating, and enough timestamp information to enforce that cutoff. Fact-check pages, their domains, and verdict-bearing post-cutoff documents are excluded from the primary evidence pool.

The first authenticated candidate snapshot was frozen at 2026-10-24 07:53:33 UTC. It contains 32,898 deduplicated ClaimReview candidates and 440 preserved API response pages from seven configured publisher filters. This number is a screening-pool count rather than a benchmark size: eligibility, United States relevance, image identity, rating, event membership, and decision-time metadata still require audit. PolitiFact, AFP, and Snopes reached the preregistered 100-page collection cap, so this snapshot is not interpreted as a census of their ClaimReview records.

Evidence timestamps have three confidence levels. \emph{Confirmed} items have consistent publication metadata and an independent pre-cutoff observation. \emph{Observed} items have a pre-cutoff archive or crawl time but incomplete publication metadata. \emph{Uncertain} items rely on self-reported metadata or were first observed after the cutoff. The primary analysis uses confirmed and observed items only; uncertain items enter a sensitivity analysis. A post-cutoff archive does not establish pre-cutoff availability.

\subsection{Label Policy}

Only explicit site ratings are normalized through a frozen, site-specific mapping. True or correct maps to supported; false, fake, or miscaptioned maps to refuted; and unproven, unsubstantiated, or no evidence maps to insufficient evidence. Mixed, partially true, satire, opinion, and ambiguous ratings are excluded. The labels come from professional fact-checkers, and automatically extracted evidence links remain silver evidence.

The evidence-sufficiency target is audited separately. It is not inferred solely from the professional verdict: a later supported or refuted verdict does not prove that enough time-valid evidence existed at an earlier cutoff. A 200-case pilot receives gold sufficiency judgments under a written protocol. Any model trained with automatically derived sufficiency targets reports those labels as silver and evaluates them separately.

\subsection{Temporal Partitioning and Release}

Training covers 2020--2023, validation uses 2024, and final testing uses 2025 plus a partial-2026 snapshot ending on 2026-10-24. The 2025 and 2026 results are reported separately and jointly. Events, near-duplicate claims, and images with matching perceptual hashes remain in one partition. The release package contains URLs, hashes, timestamp confidence, normalized labels, graph annotations, retrieval queries, and split manifests. It does not redistribute copyrighted articles or images.

Figure~\ref{fig:temporal-protocol} shows the order of operations. Group identities are fixed before the year-based allocation. Training data are used to fit model parameters and preprocessing; validation data are reserved for model selection, calibration, and abstention thresholds; the final period is not used for any of those choices.

\begin{figure}[t]
\centering
\resizebox{\columnwidth}{!}{\begin{tikzpicture}[
  font=\sffamily\footnotesize,
  phase/.style={draw, very thick, rounded corners=1.5pt, minimum height=11mm, align=center},
  guard/.style={draw=tpggray, rounded corners=1.5pt, fill=black!3, align=center, inner sep=2.2mm}
]
  \node[phase, draw=tpgblue, fill=tpgblue!9, text width=32mm] (train) {
    \textbf{Train}\\2020--2023
  };
  \node[phase, draw=tpgteal, fill=tpgteal!9, text width=20mm, right=1.5mm of train] (val) {
    \textbf{Validate}\\2024
  };
  \node[phase, draw=tpgorange, fill=tpgorange!9, text width=27mm, right=1.5mm of val] (test) {
    \textbf{Final test}\\2025 + partial 2026
  };

  \draw[-{Stealth[length=2mm]}, very thick, draw=tpggray]
    ([yshift=8mm]train.west) -- ([yshift=8mm]test.east)
    node[midway, above=0.7mm, font=\sffamily\scriptsize] {chronological allocation};

  \node[guard, text width=84mm, below=4mm of val] (lock) {
    \textbf{Group lock before allocation:}
    event identity $\mid$ near-duplicate claim $\mid$ perceptually identical image
  };
  \node[font=\sffamily\scriptsize, text=tpgorange, above=0.8mm of test] {snapshot: 2026-10-24};
  \draw[-{Stealth[length=1.8mm]}, semithick, draw=tpggray] (lock.north -| train) -- (train.south);
  \draw[-{Stealth[length=1.8mm]}, semithick, draw=tpggray] (lock.north -| val) -- (val.south);
  \draw[-{Stealth[length=1.8mm]}, semithick, draw=tpggray] (lock.north -| test) -- (test.south);

\end{tikzpicture}
}
\caption{Primary temporal split and leakage controls. A locked event, near-duplicate claim, or reused image cannot cross partitions.}
\label{fig:temporal-protocol}
\end{figure}

This study treats \datasetname{} as a benchmark only if the filtered corpus contains at least 2,000 cases, 300 cases per label, 50 test events, and two independent evidence clusters for at least 70\% of retained claims. If the collection misses any threshold, it becomes a temporal and provenance stress test. The threshold is fixed before inspecting the full collection.

\IfFileExists{tables/generated_dataset_summary.tex}{%
\input{tables/generated_dataset_summary}%
}{}

